\subsubsection{Relatórios em PDF}
Pseudo-códigos das funções de relatório, revisados: os nomes de campo passam a corresponder exatamente ao modelo (\textit{data\_devolucao\_efetuada}, não \textit{data\_devolucao}; filtro por \textit{id\_almoxarifado} denormalizado em \textit{Emprestimo\_Reagente}, não por um inexistente \textit{id\_local\_saida}), e \textit{gerarRelatorioAlmoxarifado} passa a checar o escopo do gestor, não apenas o papel.

\begin{lstlisting}[language=TypeScript, title={Configuração Base e Relatório de Almoxarifado}]
import PDFDocument from "pdfkit-table";

// PDF-026: Diferenciação Nominal vs Físico e Extravio
// 1. O relatório obedece à seguinte semântica de dados (PDF-026):
//    medida_utilizada -> consumo histórico
//    peso_atual -> última medição bruta
//    peso_frasco_vazio -> tara quando conhecida
//    conteudo_nominal -> declaração original do fabricante
//    saldo atual -> derivado apenas quando houver dados físicos suficientes.
// 2. Frascos com estado EXTRAVIADO não devem ser contados como aptos para retirada. O predicado exige estado_fisico_frasco IN {FECHADO, ABERTO}.
//    Eles devem aparecer segregados em seções de perdas/extravios.

interface FiltrosAlmoxarifado {
  idAlmoxarifado: string;
  mes: number;
  ano: number;
}

export const gerarRelatorioAlmoxarifado = onCall(async (request) => {
  const { idAlmoxarifado, mes, ano } = request.data as FiltrosAlmoxarifado;
  await validarPermissao(request, ["Chefe_Geral", "Gestor_Almoxarifado"]);
  await validarGestorDoAlmoxarifado(request.auth!.uid, request.auth!.token, idAlmoxarifado);

  const dataInicio = new Date(ano, mes - 1, 1);
  const dataFim = new Date(`${ano}-${String(mes).padStart(2, "0")}-${new Date(ano, mes, 0).getDate()}T23:59:59.999-03:00`);

  const devolucoesSnapshot = await admin.firestore().collection("Emprestimo_Reagente")
    .where("id_almoxarifado", "==", idAlmoxarifado)
    .where("data_devolucao_efetuada", ">=", dataInicio)
    .where("data_devolucao_efetuada", "<=", dataFim)
    .get();

  const consumoPorUnidade = devolucoesSnapshot.docs.reduce((acc, doc) => {
    const d = doc.data();
    if (d.unidade_medida_utilizada === "ml") acc.ml += d.medida_utilizada || 0;
    else if (d.unidade_medida_utilizada === "g") acc.g += d.medida_utilizada || 0;
    return acc;
  }, { g: 0, ml: 0 });

  const historicoSnapshot = await admin.firestore().collection("Historico_Frasco_Reagente")
    .where("id_almoxarifado", "==", idAlmoxarifado)
    .where("timestamp", ">=", dataInicio)
    .where("timestamp", "<=", dataFim)
    .orderBy("timestamp", "asc").get();

  const doc = new PDFDocument({ margin: 30, size: "A4" });
  const buffer = await buildPdfBuffer(doc, historicoSnapshot.docs, consumoPorUnidade);

  return { pdfBase64: buffer.toString("base64") };
});
\end{lstlisting}

\begin{lstlisting}[language=TypeScript, title={Relatório Mensal de um Prédio (Filtros Compostos)}]
interface FiltrosPredio {
  predio: string;
  mes: number;
  ano: number;
  andar?: string;
  sala?: string;
  estadoConservacao?: string;
  status?: string;
}

export const gerarRelatorioBensPredio = onCall(async (request) => {
  const filtros = request.data as FiltrosPredio;
  await validarPermissao(request, ["Chefe_Geral", "Gestor_Bens_Patrimoniais"]);

  let query = admin.firestore().collection("Bem_Patrimonial").where("predio", "==", filtros.predio);
  if (filtros.andar) query = query.where("andar", "==", filtros.andar);
  if (filtros.sala) query = query.where("sala", "==", filtros.sala);
  if (filtros.estadoConservacao) query = query.where("estado_conservacao", "==", filtros.estadoConservacao);
  if (filtros.status) query = query.where("status", "==", filtros.status);

  const bensSnapshot = await query.get();

  const doc = new PDFDocument({ size: "A4" });
  const buffer = await buildPredioPdfBuffer(doc, bensSnapshot.docs, filtros);

  return { pdfBase64: buffer.toString("base64") };
});
\end{lstlisting}

\begin{lstlisting}[language=TypeScript, title={Relatório Geral e Período Personalizado}]
interface FiltrosGeralEPersonalizado {
  dataInicio: string;
  dataFim: string;
  entidade: "Bens_Patrimoniais" | "Reagentes";
  prediosSelecionados?: string[];
}

export const gerarRelatorioPersonalizado = onCall(async (request) => {
  const { dataInicio, dataFim, entidade, prediosSelecionados } = request.data as FiltrosGeralEPersonalizado;
  await validarPermissao(request, ["Chefe_Geral", "Gestor_Bens_Patrimoniais", "Gestor_Almoxarifado"]);

  const dataIniObj = new Date(dataInicio);
  const dataFimObj = new Date(dataFim);
  let snapshot;

  if (entidade === "Bens_Patrimoniais") {
    let query = admin.firestore().collectionGroup("Historico") // somente caminhos Bem_Patrimonial/{id}/Historico
      .where("timestamp", ">=", dataIniObj)
      .where("timestamp", "<=", dataFimObj);
    if (prediosSelecionados && prediosSelecionados.length > 0) {
      query = query.where("predio", "in", prediosSelecionados);
    }
    snapshot = await query.get();
  } else {
    snapshot = await admin.firestore().collection("Emprestimo_Reagente")
      .where("data_devolucao_efetuada", ">=", dataIniObj)
      .where("data_devolucao_efetuada", "<=", dataFimObj)
      .get();
  }

  const doc = new PDFDocument({ layout: "landscape", size: "A4" });
  const buffer = await buildPersonalizadoPdfBuffer(doc, snapshot.docs, entidade);

  return { pdfBase64: buffer.toString("base64") };
});
\end{lstlisting}

\begin{lstlisting}[language=TypeScript, title={Módulo de Etiquetas em PDF: Virgens e Reimpressão}]
// ============================================================================
// GERACAO DE ETIQUETAS VIRGENS (NOVOS FRASCOS) - ABA 1
// ============================================================================
interface DadosEtiquetasVirgens {
  codigoInicial: number;
  codigoFinal: number;
  startRow?: number; // 1 a 10 (Grid A4)
  startCol?: number; // 1 a 3
}

export const gerarPdfEtiquetasVirgens = onCall(async (request) => {
  await validarPermissao(request, ["Chefe_Geral", "Gestor_Almoxarifado"]);
  const dados = request.data as DadosEtiquetasVirgens;

  const total = dados.codigoFinal - dados.codigoInicial + 1;
  if (total <= 0 || total > 50) {
    throw new HttpsError("invalid-argument", "O lote deve conter entre 1 e 50 etiquetas por impressão.");
  }

  // Audita a sessao de geracao de etiquetas virgens (Secao 4.13)
  await admin.firestore().collection("Impressao_Etiqueta_Frasco").add({
    gerado_em: admin.firestore.FieldValue.serverTimestamp(),
    gerado_por: request.auth!.uid,
    codigo_inicial: dados.codigoInicial,
    codigo_final: dados.codigoFinal,
  });

  const codigos = Array.from({ length: total }, (_, i) => `LCQUI-${dados.codigoInicial + i}`);
  const buffer = await gerarPdfEtiquetasNovosFrascos(codigos, dados.startRow || 1, dados.startCol || 1);

  const fileName = `etiquetas/virgens_${dados.codigoInicial}_a_${dados.codigoFinal}_${Date.now()}.pdf`;
  const fileRef = admin.storage().bucket().file(fileName);
  await fileRef.save(buffer, { contentType: "application/pdf" });
  const [url] = await fileRef.getSignedUrl({ action: "read", expires: Date.now() + 3600000 });

  return { url };
});

// ============================================================================
// REIMPRESSAO E FICHA DE CONFERENCIA (FRASCOS JA CADASTRADOS) - ABA 2
// ============================================================================
interface DadosReimpressao {
  frascoIds: string[]; // Máximo 10 frascos
}

export const gerarPdfReimpressaoFrascos = onCall(async (request) => {
  await validarPermissao(request, ["Chefe_Geral", "Gestor_Almoxarifado"]);
  const { frascoIds } = request.data as DadosReimpressao;

  // Regra 7.2.10: Teto rigido de seguranca quimica
  if (!frascoIds || frascoIds.length === 0 || frascoIds.length > 10) {
    throw new HttpsError("invalid-argument", "Selecione entre 1 e 10 frascos por sessao de reimpressao.");
  }

  // 1. Busca dados completos, historico recente e ultimo emprestimo dos frascos
  const dadosFrascos = await buscarDadosCompletosParaConferencia(frascoIds);

  // 2. Registra auditoria individual de 2a via para cada frasco (Secao 4.41)
  const batch = admin.firestore().batch();
  for (const frascoId of frascoIds) {
    const auditRef = admin.firestore().collection("Registro_de_Auditoria").doc();
    batch.set(auditRef, {
      id_usuario: request.auth!.uid,
      acao: "REIMPRESSAO_ETIQUETA",
      tipo_entidade_sofre_acao: "FRASCO_REAGENTE",
      id_do_objeto_da_entidade: frascoId,
      acao_feita_em: admin.firestore.FieldValue.serverTimestamp(),
      metadata: { motivo: "segunda_via_conferencia" },
    });
  }
  await batch.commit();

  // 3. Gera PDF com Ficha de Conferencia + 1 Etiqueta Centralizada por frasco
  const buffer = await gerarPdfFrascosCadastrados(dadosFrascos);

  const fileName = `etiquetas/reimpressao_${Date.now()}.pdf`;
  const fileRef = admin.storage().bucket().file(fileName);
  await fileRef.save(buffer, { contentType: "application/pdf" });
  const [url] = await fileRef.getSignedUrl({ action: "read", expires: Date.now() + 3600000 });

  return { url };
});
\end{lstlisting}



